%% ma: L10-B01-0039
%% lop: 10
%% chuong: 1
%% bai: 1
%% dang: 10.1.7
%% loai: TN4
%% muc: VD
%% dap_an: "-"
%% nguon: "Ảnh giáo viên gửi 10/10/2026 (ThS. Nguyễn Hoàng Việt, Chương 1 Mệnh đề), B-Đề 2, Phần 1 Câu 10"
%% kien_thuc: "Bài 1 (mệnh đề; suy luận logic, tổ hợp)"
%% trang_thai: cach-ly
%% ly_do: "Theo đề, n lớn nhất là 5 (n=6 cần phủ 15 cặp tầng nhưng 4 thang máy chỉ phủ tối đa 12 cặp), không có trong bốn phương án 6, 7, 8, 9. Gợi ý: đề thiếu hoặc nhầm điều kiện, cần đối chiếu sách gốc"
%% ghi_chu: "Đã kiểm bằng Python: n=5 là giá trị lớn nhất thỏa điều kiện (n=3,4 và 6,7 không thỏa)"
\begin{ex}
Một tòa nhà có $n$ tầng, các tầng được đánh số từ $1$ đến $n$ theo thứ tự từ dưới lên. Có $4$ thang máy đang ở tầng $1$. Biết rằng mỗi thang máy có thể dừng ở đúng $3$ tầng và $3$ tầng này không là $3$ số nguyên liên tiếp và với hai tầng bất kỳ của tòa nhà luôn có một thang máy dừng được ở cả hai tầng này. Hỏi giá trị lớn nhất của $n$ là bao nhiêu?
\choice{$6$}{$7$}{$8$}{$9$}
\loigiai{Mỗi thang máy phủ được $\binom32=3$ cặp tầng, bốn thang máy phủ tối đa $12$ cặp, trong khi $n$ tầng có $\binom n2$ cặp; cần $\binom n2\le12$ nên $n\le5$. Với $n=5$ xếp được (kiểm bằng Python), $n=6$ thì không.}
\end{ex}
